%!TEX root = forallxyyc.tex
% Capa

\thispagestyle{empty}

\begin{textblock*}{\stockwidth}(0in,3.5in)
  \noindent\includegraphics{../assets/forallx-coverart-screen.pdf}
\end{textblock*}

% criar a capa
\begin{textblock*}{0in}(0in,0in)
  \noindent\hfill
  \begin{minipage}[b][\stockheight][s]{.9\stockwidth}
    \begin{raggedleft}
      \vspace*{1.7cm}
      \hfill
      \sffamily\fontsize{66pt}{0pt}\selectfont
      \color{lyallpink}
      \textbf{forall}%
      \fontsize{80pt}{0pt}\selectfont\rmfamily\textit{\textbf{x}\/}
                    
      \vskip.5cm
          
      \sffamily\fontsize{24pt}{0pt}\selectfont
      \color{black}
      \textbf{\uppercase{Calgary}}

      \vspace*{1cm}
      \color{black}
      \sffamily
      \fontsize{20pt}{22pt}\selectfont
      \textbf{Soluções para\\ Exercícios selecionados}

      \vfill
      \fontsize{11pt}{14pt}\selectfont
      \color{white}
      \textbf{P.~D. Magnus}\\
      \textbf{Tim Button}\\
      \textbf{Richard Zach}\\
      \textit{com contribuições de}\\
      \textbf{Robert Trueman}\\
      \textbf{J.~Robert Loftis}\\
      \textbf{Aaron Thomas-Bolduc}\\ \par
      \vspace{1cm}
      \textbf{\forallxversion}\par
      \vspace{1cm}
    \end{raggedleft}
  \end{minipage}
  \hspace*{1cm}
\end{textblock*}
\ 
\newpage
\color{black}
\setlength{\barlength}{0pt} % nenhuma barra ao lado dos títulos de capítulo no e-book

\noindent Este caderno baseia-se no caderno de soluções de \href{https://www.homepages.ucl.ac.uk/~uctytbu/OERs.html}{\forallx: \emph{Cambridge}}, de\\[2ex]
\href{https://www.homepages.ucl.ac.uk/~uctytbu/}{Tim Button}\\
\emph{University of Cambridge}\\[2ex]
utilizado sob uma licença \href{https://creativecommons.org/licenses/by/4.0/}{CC BY 4.0}, que por sua vez se baseia em \href{https://www.fecundity.com/logic/}{\forallx}, de\\[2ex]
\href{https://www.fecundity.com/job/}{P.D.\ Magnus}\\
\emph{University at Albany, State University of New York}\\[2ex]
utilizado sob uma licença \href{https://creativecommons.org/licenses/by/4.0/}{CC BY 4.0}, e foi remixado, revisado e ampliado por\\[2ex] {Aaron
Thomas-Bolduc \& \href{https://richardzach.org/}{Richard Zach}}\\
\emph{University of Calgary}
\\
\\
Ele inclui material adicional de \forallx{} por P.D. Magnus, utilizado sob uma licença \href{https://creativecommons.org/licenses/by/4.0/}{CC BY 4.0}, e de
\href{https://github.com/rob-helpy-chalk/openintroduction}{\forallx: \emph{Lorain
County Remix}},
de \href{https://sites.google.com/site/cathalwoods/}{Cathal Woods} e
J. Robert Loftis, e de
\href{http://www.rtrueman.com/uploads/7/0/3/2/70324387/modal_logic_primer.pdf}{\emph{A Modal Logic Primer}}, de \href{http://www.rtrueman.com/}{Robert
Trueman}, utilizado com permissão.

\bigskip

\noindent Esta obra está licenciada sob uma licença \href{https://creativecommons.org/licenses/by/4.0/}{Creative Commons Atribuição 4.0}.
Você pode copiar e redistribuir o material em qualquer meio ou formato, bem como remixá-lo, transformá-lo e criar obras derivadas para qualquer finalidade, inclusive comercialmente, nos termos seguintes:
\begin{itemize}
\item Você deve atribuir o devido crédito, fornecer um link para a licença e indicar se foram feitas alterações. Isso pode ser feito de qualquer maneira razoável, mas não de modo que sugira que o licenciador apoia você ou o uso que você faz da obra.
\item Você não pode aplicar termos jurídicos nem medidas tecnológicas que restrinjam legalmente outras pessoas de fazer qualquer coisa que a licença permita.
\end{itemize}
O código-fonte \LaTeX{} deste livro está disponível no
\href{https://github.com/rzach/forallx-yyc/}{GitHub} e um PDF está disponível em
\href{https://forallx.openlogicproject.org}{forallx.openlogicproject.org}.
Esta versão é a revisão \gitAbbrevHash{} (\gitCommitterDate).

\bigskip
\noindent A preparação deste livro-texto foi possível graças a uma bolsa do \href{https://taylorinstitute.ucalgary.ca/}{Taylor Institute for Teaching and Learning}.

\bigskip
\noindent
\href{https://taylorinstitute.ucalgary.ca/}{\includegraphics[width=8cm]{../assets/ti-color}}
